\documentclass[
 %reprint,
 aps, pra,
 amsmath,amssymb,
 11pt,
 final,
% notitlepage,
tightenlines,
 twoside,
 twocolumn,
 nofloats,
% nobibnotes,
nofootinbib,
 superscriptaddress,
%noshowpacs,
showkeys,
showkeywords,
%centertags
 ]
{revtex4-2}

\usepackage[T1]{fontenc}
\usepackage[utf8x]{inputenc}
\usepackage[english]{babel}
\usepackage{graphicx}% Include figure files
\usepackage{dcolumn}% Align table columns on decimal point
\usepackage{bm}% bold math
%\usepackage{longtable}


\input{maik.rty}


%\renewcommand{\rmdefault}{lh}

\setcitestyle{authoryear,round}
\setlength{\bibhang}{1.5em}





%Подключение некоторых команд из файла sao_cmd_author.tex
\input{sao_cmd_author.tex}


% Начало документа

\begin{document}

\selectlanguage{english}

\keywords{stars: double or binary---stars: individual: ADS\,48}

%\ydk{}
%\titlerunning{}
%\authorrunning{}
%\toctitle{}
%\tocauthor{}

\title{The Title of the Paper}

\author{\firstname{I.}~\surname{Author1}}
 \email{example@list.ru}
 \affiliation{Main (Pulkovo) Astronomical Observatory
of the RAS, St. Petersburg, 196140 Russia}

\author{\firstname{I.}~\surname{Author2}}
 \affiliation{\saoname}


\begin{abstract}
We study the nearest outskirts ($R<3R_{200c}$) of 40 groups and
clusters of galaxies of the local Universe ($0.02<z<0.045$ and
300~km~s$^{-1}<\sigma<950$~km~s$^{-1}$). Using the SDSS\,DR10
catalog data, we determined the stellar mass of galaxy clusters
corresponding to $K_s$-luminosity (which we determined earlier
based on the 2MASX catalog data) $(M_*/M_{\odot}) \propto
(L_K/L_{\odot})^{1.010\pm0.004}$ (\mbox{$M_K<-21\fm5$},
$R<R_{200c}$). We also found the dependence of the galaxy cluster
stellar mass on halo mass: \mbox{$(M_*/M_{\odot}) \propto
(M_{200c}/M_{\odot})^{0.77\pm0.01}$}\ldots
\end{abstract}

\maketitle

\section{INTRODUCTION}

The multiple star system ADS\,48, discovered by Otto Struve in 1876,
has been repeatedly investigated by various authors~\citep[see,
for example,][]{G-L, Hopman}, but their attention was mainly attracted
by the inner pair AB. According to the identification in the \citet{WDS}
catalog, the three stars---A, B and F---are physically connected
(by common parallax and proper motions). For the inner AB pair,
there has been a long series of observations since its discovery,
and the F component has not been observed for almost a century.


\section{FIRST SECTION}

Here comes some math:
\begin{eqnarray}
\rho&=&\sqrt{x^2+y^2}~,~~~\, \qquad \theta\;=\;\arctan\dfrac{x}{y} \nonumber \\
\mu&=&\sqrt{\mu_{x}'^2+\mu_{y}'^2}~, \qquad
\psi\;=\;\arctan\dfrac{\mu_{x}'}{\mu_{y}'}. \nonumber
\end{eqnarray}
Here
\begin{eqnarray}
x&=&(\alpha_{\rm B}-\alpha_{\rm A})\cos{\delta}\times3600,\nonumber\\
y&=&(\delta_{\rm B}-\delta_{\rm A})\times3600,\nonumber\\
\delta&=&(\delta_{\rm A}+\delta_{\rm B})/2,\nonumber\\
\mu_{x}'&=&\mu_{x{\rm B}}-\mu_{x{\rm A}},\nonumber
\qquad\mu_{y}'=\mu_{y{\rm B}}-\mu_{y{\rm A}}\nonumber.
\end{eqnarray}

%1
\begin{table*}
\caption {Two column table} \label{pvd}
\medskip
\begin{tabular}{c|c|c|c|c|c}
\hline
pair                           & AB             & AB                    & AB--F             & AB--F              & AB--F \\
\hline
\multirow{2}{*}{Instrument}    & $26''$, CCD    & \multirow{2}{*}{GAIA} & $26\arcsec$, photo & $26\arcsec$, photo  & \multirow{2}{*}{GAIA} \\
                               & individual &                       & individual    & smoothed         &    \\
\cline{1-1}
Parameters                      &                &                       &                   &                    &    \\
\hline
$T_1-T_2$                        & $2003$--$2012$ & --           & $1968$--$1995$ & 1971--1992  & -- \\
$T_0$                            & $2008.6$       & $2015.5$     & $1981.5$       & 1981.5      & $2015.5$ \\
  n                              & $48$           & --           & $115$          & 30          & -- \\
 \hline
$\rho$, arcsec.                & $6.0534$       & $6.00768^*$  & $327.3322$     & $327.3339$  & $327.1754$ \\
                                 & $\pm0.0012$    & $\pm0.00008$ & $\pm0.0023$    & $\pm0.0010$ & $\pm0.0002$ \\
\hline
$\theta_{2000}$, degr.         & $185.3604$     & $188.2084$   & $254.2942$     & $254.2943$  & $254.25739$ \\
                                 & $\pm0.0059$    & $\pm0.0010$  & $\pm0.0017$    & $\pm0.0005$ & $\pm0.00001$ \\
\hline
$\mu$, mas/year                   & $43.1$         & $44.94$      & $4.3$          & $3.9$       & $5.4$ \\
                                 & $\pm0.3$       & $\pm0.18$    & $\pm0.5$       & $\pm0.2$    & $\pm0.1$ \\
\hline
$\psi_{2000}$, degr.             & $283.09$       & $288.06$     & $86.4$         & $73.2$      & $37.8$ \\
                                 & $\pm0.7$       & $\pm0.16$    & $\pm19.6$      & $\pm7.4$    & $\pm1.1$ \\
\hline
$\dot{\rho}$, mas/year            & $-6.4$         & $-7.7$       & $-4.2$         & $-3.9$      & $-4.4$ \\
                                 & $\pm0.5$       & $\pm0.1$     & $\pm0.4$       & $\pm0.2$    & $\pm0.1$ \\
\hline $\dot{\theta}$, degr./year & $0.4034$       & $0.4202$     & $-0.0002$      & $-0.0000$   & $-0.00056$ \\
                                 & $\pm0.0024$    & $\pm0.0017$  & $\pm0.0003$    & $\pm0.0001$ & $\pm0.00001$ \\
\hline
\multicolumn{6}{l}{\footnotesize {Here $n$ is the number
of individual or smoothed
observations,}} \\[-5pt]
\multicolumn{6}{l}{\footnotesize {$^*$---the value $\rho$
is given adjusted for Gaia--CCD$=+0\farcs03$.}}\\
\end{tabular}
\end{table*}

\begin{figure*}
\includegraphics[scale=0.3]{Kiyaeva_fig1a.eps}
\includegraphics[scale=0.3]{Kiyaeva_fig1b.eps}
\caption{Two column figure.}
 \label{G-CCD}
\end{figure*}

Table~\ref{pvd} shows AMPs calculated from observations of Gaia\,DR2,
and long-term series of observations of the Pulkovo 26-inch refractor.
For the AB pair, we compare only with the CCD observations of 2003--2012.
For this pair, we found a systematic discrepancy in $\rho$,
which is clearly visible in Fig.~\ref{G-CCD}.



\begin{figure}
\includegraphics[scale=0.3]{Kiyaeva_fig2a.eps}
\includegraphics[scale=0.3]{Kiyaeva_fig2b.eps}
\caption{One column figure.}
 \label{G-photo}
\end{figure}

\section{ANOTHER SECTION}

The text of the Section.



%\begin{figure*}
%\includegraphics[scale=0.33]{Kiyaeva_fig3a.eps}
%\includegraphics[scale=0.33]{Kiyaeva_fig3b.eps}
%\includegraphics[scale=0.35]{Kiyaeva_fig3c.eps}
%\caption{Three-part figure.}
%\label{figAB}
%\end{figure*}


\section{THIRD SECTION}

The motion of the outer pair is in the direction $\rho$, and we can
definitely state that for all orbits of the family,
the inclination of the orbit $i\approx90\degr$,
and the longitude of the ascending node $\Omega\approx\theta-180\degr$.
Therefore, you can calculate the angle between the planes of
the outer and inner orbits. As a result, we get that the planes
of the orbits are non-planar. 


Equations:
\begin{equation}
\label{one} 
v_1=\sqrt{\dfrac{4\pi^2{m_2}^2}{r(m_1+m_2)}}.\nonumber
\end{equation}
Assuming $m_2<<m_1$, we get
\begin{equation}
\label{two} 
m_2=v_1 \times\sqrt{\dfrac{m_1}{4\pi^2}r},\nonumber
\end{equation}

\begin{equation}
\label{three}
\overrightarrow{v_1}=f\times(\overrightarrow{\mu_{\rm
G}}-\overrightarrow{\mu_{\rm ph}})/p_t,\nonumber
\end{equation}
where $\overrightarrow{\mu_{\rm ph}}=(\mu_{\rm ph} \sin\,\psi_{\rm
ph},\mu_{\rm ph} \cos\,\psi_{\rm ph})$is the average orbital motion
obtained from a long series of photographic observations;
$\overrightarrow{\mu_{\rm G}}$ is the instantaneous orbital motion
determined by Gaia observation; \mbox{$p_t=87$}~mas is the parallax; $f$ is the coefficient of transition
from the relative velocity of the orbital motion to the velocity
relative to the center of mass of the hierarchical triple system,
which is motionless. If component F has a satellite,
then $$f_{\rm F}=M_{\rm A+B}/M_{\rm A+B+F}.$$ If the center of mass
of the AB system oscillates, then
\begin{eqnarray}
 f_{\rm C}&=&M_{\rm F}/M_{\rm A+B+F},\nonumber\\
 f_{\rm A}&=&f_{\rm C}\,(M_{\rm A+B}/M_{\rm A}),\nonumber\\
 f_{\rm B}&=&f_{\rm C}\,(M_{\rm A+B}/M_{\rm B}).\nonumber
\end{eqnarray}

If we use the values of $\mu_{ph}$ according to the smoothed series,
then
\begin{eqnarray}
 m_{2,{\rm F}}/\sqrt{r}&=&0.0030\pm0.0006~M_\odot,\nonumber\\
 m_{2,{\rm A}}/\sqrt{r}&=&m_{2,{\rm B}}/\sqrt{r}=0.0027\pm0.0006~M_\odot;\nonumber
\end{eqnarray}
if we use the values of $\mu_{ph}$ according to individual observations, then
\begin{eqnarray}
 m_{2,{\rm F}}/\sqrt{r}&=&0.0039\pm0.0017~M_\odot,\nonumber\\
 m_{2,{\rm A}}/\sqrt{r}&=&m_{2,{\rm B}}/\sqrt{r}=0.0035\pm0.0015~M_\odot.\nonumber
\end{eqnarray}

%4
\begin{table}
\caption {One column table.} \label{orbF}
\medskip
\begin{tabular}{c|c|c|c|c}
\hline
\multirow{2}{*}{Parameters} & \multicolumn{4}{c}{Orbits}  \\
\cline{2-5}
 & 1 & 2 & 3 & 4 \\
\hline
$a_{\rm ph}$,~mas          & 15.0    & 14.3    & 8.2     & 4.0 \\
$P$,~year                   & 11.0    & 11.04   & 9.52    & 10.97 \\
$e$                        & 0.2     & 0.24    & 0.53    & 0.3 \\
$i$,~degr.                 & 97.0    & 96.3    & 179.98  & 44 \\
$\omega$,~degr.            & 235.0   & 258.6   & 79.8    & 56.2 \\
$\Omega$,~degr.            & 147.2   & 143.2   & 12.0    & 217.1 \\
$T$,~year                   & 1980.0  & 1980.56 & 1988.15 & 1982.8 \\
$V_{r\gamma}$,~m\,c$^{-1}$ & --      & --      & $-0.7$  & -- \\
\hline
$p_t$,~mas                 & 87.0    & 87.0    & 87.0    & 86.9\\
$M_1,~M_\odot$             & 0.5     & 0.5     & 0.5     & 0.65 \\
$a_1$,~a.u.                & 0.17    & 0.16    & 0.094   & 0.046 \\
$M_2,~M_\odot$             & 0.023   & 0.022   & 0.013   & 0.007 \\
$a_2$,~a.u.                & 3.82    & 3.82    & 3.50    & 4.28 \\
\hline
$\sigma_x$,~mas            & 2.2     & 2.0     & 3.6     & 5.1 \\
$\sigma_y$,~mas            & 12.3    & 12.0    & 13.5    & 5.4 \\
$\sigma_{V_r}$,~mas/year    & --      & --      & 0.078   & -- \\
\hline
\end{tabular}
\end{table}

The system of equations is solved:
\begin{equation}
\label{four} x(t)=x_0+\dot{x}(t-t_0)+BX_{\varphi}+GY_{\varphi},
\end{equation}
\begin{equation}
\label{five} y(t)=y_0+\dot{y}(t-t_0)+AX_{\varphi}+FY_{\varphi},
\end{equation}
where $x=\rho\sin\theta$, $y=\rho\cos\theta$; phase
$\varphi=(t-t_0)/P$; $X_{\varphi}=\cos\,(E_{\varphi})-e$,
$Y_{\varphi}=\sqrt{1-e^2}\sin\,(E_{\varphi})$are orbital
coordinates corresponding to dynamic orbital elements $P$, $T$
and $e$; $x_0$ and $y_0$ are coordinates of the center of mass
at the moment $t_0$; $A$, $B$, $F$ and $G$ are the Thiele-Innes
elements, from which we obtain the geometric elements of the orbit
($a$, $i$, $\omega$, $\Omega$). In Table~\ref{orbF}, this orbit
is represented by number~2.


\section{CONCLUSION}

This paper demonstrates the possibility\ldots


\begin{acknowledgments}
The authors are grateful\ldots
\end{acknowledgments}

\section*{FUNDING}
This work was supported by the\ldots

\section*{CONFLICT OF INTEREST}
The authors declare no conflicts of interest.

\bibliographystyle{aspb1}
\bibliography{Kiyaeva}

%The bibliography without bibtex:

%\begin{thebibliography}{25}
%\providecommand{\natexlab}[1]{#1}
%
%\bibitem[{Gaia Collaboration} et~al.(2018)]{Gaia1}
% A.~G.~A. {Brown} et~al. ({Gaia Collab.}), \aap
%  \textbf{616}, A1 (2018).
%
%\bibitem[{Butler} et~al.(2017)]{Butler2017}
%R.~P.~{Butler}, S.~S.~{Vogt}, G.~{Laughlin}, et~al., \aj
%\textbf{153}~(5), 208  (2017).
%
%\bibitem[{Gontcharov} et~al.(2000)]{Gonch2000}
%G.~A.~{Gontcharov}, A.~A.~{Andronova}, and O.~A.~{Titov}, \aap
%\textbf{355},  1164 (2000).
%
%\bibitem[{Gontcharov} and {Kiyaeva}(2002)]{Gonch2002}
%G.~A.~{Gontcharov} and O.~V.~{Kiyaeva}, \aap \textbf{391}, 647
%(2002).
%
%\bibitem[{Gontcharov} and {Kiyaeva}(2010)]{Gonch2010}
%G.~A.~{Gontcharov} and O.~V.~{Kiyaeva}, \na \textbf{15}~(3), 324
%(2010).
%
%\bibitem[{G{\"u}ntzel-Lingner}(1955)]{G-L}
%U.~{G{\"u}ntzel-Lingner}, \an \textbf{282}~(4), 183
%  (1955).
%
%\bibitem[{Hartkopf} and {Mason}(2015)]{Hartkopf2015}
%W.~I.~{Hartkopf} and B.~D.~{Mason}, \aj \textbf{150}~(4), 136
%(2015).
%
%\bibitem[{Hartkopf} and {Mason}(2016)]{ORB6}
%W.~I.~{Hartkopf} and B.~D.~{Mason}, USNO, Washington, version 2016
%(2016).
%
%\bibitem[{Hopmann}(1964)]{Hopman}
%J.~{Hopmann}, Annalen der K.K. Sternwarte Wien \textbf{26}, 1
%(1964).
%
%\bibitem[{Izmailov}(2019)]{Izm2019}
%I.~S.~{Izmailov}, \alet \textbf{45}~(1), 30 (2019).
%
%\bibitem[{Izmailov} et~al.(2010)]{Izm2010}
%I.~S.~{Izmailov}, M.~L.~{Khovricheva}, M.~Y.~{Khovrichev}, et~al.,
%\alet \textbf{36}~(5), 349 (2010).
%
%\bibitem[{Izmailov} and {Roshchina}(2016)]{Izm2016}
%I.~S.~{Izmailov} and E.~A.~{Roshchina}, \ab \textbf{71}~(2),
%  225 (2016).
%
%\bibitem[{Kervella} et~al.(2019)]{Kerv}
%P.~{Kervella}, F.~{Arenou}, F.~{Mignard}, and F.~{Th{\'e}venin},
%\aap  \textbf{623}, A72 (2019).
%
%\bibitem[{Khovrichev} et~al.(2016)]{Chovr2015}
%M.~Y.~{Khovrichev}, A.~M.~{Kulikova}, E.~N.~{Sokov}, et~al., \alet
%  \textbf{42}~(10), 686 (2016).
%
%\bibitem[{Kiselev} and {Kiyaeva}(1980)]{PVD80}
%A.~A.~{Kiselev} and O.~V.~{Kiyaeva}, \azh \textbf{57}, 1227
%(1980).
%
%\bibitem[{Kiselev} et~al.(2014)]{Kis2014}
%A.~A.~{Kiselev}, O.~V.~{Kiyaeva}, I.~S.~{Izmailov}, et~al., \arep
%  \textbf{58}~(2), 78 (2014).
%
%\bibitem[{Kiyaeva} et~al.(2001)]{Kiya2001}
%O.~V.~{Kiyaeva}, A.~A.~{Kiselev}, E.~V.~{Polyakov}, and
%V.~B.~{Rafal'skii},  \alet \textbf{27}, 391 (2001).
%
%\bibitem[{Mason} et~al.(2013)]{Mason2013}
%B.~D.~{Mason}, W.~I.~{Hartkopf}, and H.~M.~{Hurowitz}, \aj
%\textbf{146}~(3), 56  (2013).
%
%\bibitem[{Mason} et~al.(2016)]{WDS}
%B.~D.~{Mason}, G.~L.~{Wycoff}, W.~I.~{Hartkopf}, et~al., VizieR
%Online Data  Catalog B/wds (2016).
%
%\bibitem[{Neves} et~al.(2013)]{Neves2013}
%V.~{Neves}, X.~{Bonfils}, N.~C.~{Santos}, et~al., \aap
%\textbf{551}, A36  (2013).
%
%\bibitem[{Orlov} and {Zhuchkov}(2005)]{dyst1}
%V.~V.~{Orlov} and R.~Y.~{Zhuchkov}, \arep \textbf{49}~(3), 201
%  (2005).
%
%\bibitem[{Popovi{\'c}} and {Pavlovi{\'c}}(1996)]{Pop}
%G.~M.~{Popovi{\'c}} and R.~{Pavlovi{\'c}}, Bulletin Astronomique
%de Belgrade  \textbf{153}, 57 (1996).
%
%\bibitem[{Tal-Or} et~al.(2019)]{Tal2019}
%L.~{Tal-Or}, T.~{Trifonov}, S.~{Zucker}, et~al., \mnras
%\textbf{484}~(1), L8  (2019).
%
%\bibitem[{Tokovinin}(1990)]{Tok1990}
%A.~A.~{Tokovinin}, \emph{{Catalogue of stellar radial velocities.
%Catalogue of proper motions}} (Moscow State Univ., Moscow, 1990).
%%А.~А.~Токовинин в \emph{Каталог измерений лучевых скоростей звезд.
%%Каталог собственных движений.} (Mоск.Гос.Университет,Москва,1990),
%%92.
%
%\bibitem[{Zhuchkov} and {Orlov}(2005)]{dyst2}
%R.~Y.~{Zhuchkov} and V.~V.~{Orlov}, \arep \textbf{49}~(4), 274
%  (2005).
%
%\end{thebibliography}


\end{document}